\chapter*{ЗАКЛЮЧЕНИЕ}
\addcontentsline{toc}{chapter}{ЗАКЛЮЧЕНИЕ}

В ходе выполнения первого этапа НИОКР по теме <<Разработка и тестирование
прототипа программного комплекса для моделирования и симуляции роботов с
использованием графических ускорителей>> получены следующие основные результаты:

\begin{enumerate}
  \item Разработан и реализован алгоритм загрузки и отрисовки моделей роботов в
        формате URDF: разбор описания, построение дерева звеньев с вычислением
        мировых преобразований и интерактивная визуализация средствами
        VulkanSceneGraph.
  \item Реализован модуль загрузки моделей сцен и объектов окружения в форматах
        SDF и STL с приведением к единому внутреннему представлению сцены.
  \item Разработан и программно реализован алгоритм генерации кинематической
        схемы на основе формата STEP: восстановление дерева звеньев и суставов из
        геометрии сборки по аналитическим признакам контактов поверхностей и
        экспорт модели.
  \item Разработан и реализован алгоритм симуляции, не использующий графический
        процессор: симплектическое интегрирование с подшагами, широкая фаза
        (LBVH), узкая фаза (SAT) и решатель ограничений XPBD; реализация служит
        эталоном корректности.
  \item Разработан и протестирован базовый графический интерфейс пользователя
        (WPF, архитектура MVVM) с основными панелями и сервисным слоем.
  \item Проведено предварительное тестирование прототипа на искусственных
        данных; собраны и зафиксированы эталонные показатели производительности и
        подтверждён детерминизм результатов.
  \item Разработан и программно реализован алгоритм численного интегрирования
        дифференциальных уравнений движения с помощью GPU поверх переносимого
        вычислительного слоя SYCL/AdaptiveCpp (бэкенды CUDA, HIP, CPU).
\end{enumerate}

\textbf{Оценка полноты решения поставленных задач.} Все задачи, предусмотренные
календарным планом первого этапа, решены в полном объёме. Созданы независимо
тестируемые подсистемы прототипа программного комплекса и сформирован набор
эталонных показателей. Достоверность результатов подтверждена модульным и
сквозным тестированием, сопоставлением GPU- и CPU-реализаций и контролем
детерминизма. Полученные заделы достаточны для перехода к следующему этапу,
основными задачами которого являются полный перенос физического конвейера на GPU,
оценка масштабируемости и интеграция подсистем в единый прикладной комплекс.
